\documentclass[11pt,reqno]{amsart}
\usepackage[margin=1in]{geometry}
\usepackage{amsmath,amssymb,amsthm}
\usepackage[colorlinks=true,linkcolor=blue,citecolor=blue,urlcolor=blue]{hyperref}

\theoremstyle{plain}
\newtheorem{theorem}{Theorem}
\newtheorem{proposition}[theorem]{Proposition}
\newtheorem{lemma}[theorem]{Lemma}
\newtheorem{corollary}[theorem]{Corollary}
\newtheorem{conjecture}[theorem]{Conjecture}
\theoremstyle{definition}
\newtheorem{definition}[theorem]{Definition}
\theoremstyle{remark}
\newtheorem{remark}[theorem]{Remark}

\newcommand{\dd}{\,\mathrm{d}}
\newcommand{\R}{\mathbb{R}}
\newcommand{\C}{\mathbb{C}}
\newcommand{\E}{\mathbb{E}}
\newcommand{\Z}{\mathbb{Z}}
\newcommand{\Prob}{\mathbb{P}}
\newcommand{\SAO}{\mathcal{A}_\beta}
\newcommand{\TW}{\mathrm{TW}_\beta}
\newcommand{\eps}{\varepsilon}
\newcommand{\pinst}{p_\star}
\newcommand{\sech}{\operatorname{sech}}
\DeclareMathOperator{\Var}{Var}

% ---- provenance tags ----
\newcommand{\tagPr}{{\footnotesize\textsf{[established here]}}}
\newcommand{\tagP}{{\footnotesize\textsf{[proved elsewhere, cited]}}}
\newcommand{\tagT}{{\footnotesize\textsf{[to prove]}}}
\newcommand{\tagX}{{\footnotesize\textsf{[conjectural / frontier]}}}

\begin{document}

\title[A resurgent trans-series for the Tracy--Widom-$\beta$ tail]{A resurgent trans-series for
the Tracy--Widom-$\beta$ tail:\\ two directions of resurgence from a Painlev\'e~II backbone}
\author{Solomon Williams}
\address{School of Mathematics, University of Edinburgh, Edinburgh EH9 3FD, United Kingdom}
\email{solomonjwilliams635@outlook.com}
\thanks{Draft, \today. Companion to the coupled-FHN cusp program
(\texttt{W\_ResurgentTransseries\_paper}): the $q{=}1$ (Airy/fold) case treated here is the
integrable-backbone sibling of the $q{=}2$ (Weber/cusp) law $\mathcal W_\beta$, which has no
Painlev\'e reduction and is established only numerically. Provenance is tagged throughout:
\textsf{established here} / \textsf{cited} / \textsf{to prove} / \textsf{frontier}.}
\date{\today}

\begin{abstract}
The Tracy--Widom-$\beta$ ($\TW$) right tail is known to all orders only through a non-rigorous
loop-equation derivation (Borot--Nadal), and its large-$\beta$ rate function has recently been
reduced to a Painlev\'e~II (PII) equation by the optimal-fluctuation method (Comtet--Le
Doussal--Smith); neither is a resurgent derivation, and no random-operator edge law has, to our
knowledge, been placed in a Borel-summable resurgent trans-series. We supply one, and find that the tail carries
resurgent structure in \emph{two independent directions}. \emph{(Level.)} At weak noise
$\eps=\beta^{-1/2}$ we reduce the tail to a backward-Kolmogorov ODE, carry out the weak-noise WKB,
and show by an explicit affine Lax-pair map that the fluctuation (Jacobi) operator
$\mathcal J=-\partial_x^2+6\pinst^2-2(x+a)$ \emph{is} the exact Painlev\'e~II linearization
$\mu^2(-\partial_s^2+6q^2+s)$, the instanton being the $\alpha=\tfrac12$ Airy special solution of
PII. Costin's theorem then makes the $a\to\infty$ expansion the Hastings--McLeod trans-series
(Theorem~\ref{thm:T1}); its Stokes constant is computed here to $28$ digits,
$C=\sqrt{2/(3\pi^3)}=\tfrac1\pi\sqrt{2/(3\pi)}=0.146632271193848478\ldots$, matching the published
Hastings--McLeod value and refuting the natural guess $1/(4\sqrt\pi)$ (Theorem~\ref{thm:T2}). The
$a$-dependence of the Stokes constant is purely algebraic --- a Gelfand--Yaglom zero-mode Jacobian
$\sqrt{2/(3\pi)}=\lVert\psi_0\rVert/\sqrt{2\pi}$ of the tanh-kink escape instanton
(Theorem~\ref{thm:T2}(d)). \emph{(Noise.)} The weak-noise $1/\beta$ expansion of the tail
is itself Borel summable along the noise ray, in the \emph{median} sense; we prove this elementarily
--- the reduced escape has an Airy-type integral representation, so no nonlinear-ODE theorem is
needed --- and exhibit its two sectors: a dynamical escape sector (the loop factor
$\sum_n c_n(a)\beta^{-n}$; positive real axis, median) and a Coulomb-gas sector (the $\Gamma(\beta/2)$
normalization; imaginary axis, ordinary) (Theorem~\ref{thm:T3}). A clean structural dichotomy with the cusp companion results: the cusp's
complexity is a \emph{resonance} phenomenon, $\TW$'s a \emph{trans-series} one.
\end{abstract}

\maketitle

\section{Objects and the weak-noise parameter}\label{sec:objects}

\begin{definition}[Stochastic Airy operator and $\TW$]\label{def:sao}
For $\beta>0$ let $\SAO=-\dfrac{\dd^2}{\dd x^2}+x+\dfrac{2}{\sqrt\beta}\,b'(x)$ on
$L^2(\R_+)$ with a Dirichlet condition at $0$, $b'$ a white noise; let $\Lambda_0(\beta)$ be its
lowest eigenvalue and $\TW=\mathrm{law}(-\Lambda_0(\beta))$, with right tail
$\mathcal T_\beta(a):=\Prob(\TW>a)=\Prob(\Lambda_0(\beta)<-a)$. \tagP\ \emph{(Ram\'irez--Rider--Vir\'ag
\cite{RRV}.)}
\end{definition}

Write $\eps:=\beta^{-1/2}$: large $\beta$ is the weak-noise limit $\eps\to0$, and $\eps^2=1/\beta$
is the loop-expansion (trans-series) variable.

\begin{proposition}[Rigorous anchor]\label{prop:anchor}
$\mathcal T_\beta(a)=a^{-\frac34\beta+o(1)}\exp(-\tfrac23\beta a^{3/2})$ as $a\to\infty$; the leading
rate is the instanton action $\Phi(a)=\tfrac23 a^{3/2}$ and $a^{-3\beta/4}$ is the (eikonal)
prefactor. \tagP\ \emph{(Dumaz--Vir\'ag \cite{DV}.)}
\end{proposition}

\begin{remark}[The all-orders target]\label{rem:target}
The full $1/\beta$ expansion of $\mathcal T_\beta$ exists to all orders only via the physical
loop-equation derivation of Borot--Nadal \cite{BorotNadal}, which gives the exact soft tail
\begin{equation}\label{eq:bn}
\mathcal T_\beta(a)=\frac{\Gamma(\beta/2)}{(4\beta)^{\beta/2}\,2\pi}\;a^{-3\beta/4}\,
e^{-\frac23\beta a^{3/2}}\,\exp\!\Big[\sum_{m\ge1}\tfrac\beta2 R_m(2/\beta)\,a^{-3m/2}\Big],
\end{equation}
$R_m$ explicit polynomials in $2/\beta$ --- never proved, never resummed. Producing this object as a
Borel-summable resurgent trans-series is the goal. \emph{(cited; non-rigorous, unresummed.)}
\end{remark}

\section{Two directions of resurgence}\label{sec:two}

The weak-noise WKB of \S\ref{sec:T1proof} generates a loop expansion
$F(a;\beta)=\sum_{n\ge0}c_n(a)\beta^{-n}$ (formal), the fluctuation factor multiplying the leading
$e^{-\beta\Phi(a)}$, $\Phi(a)=\tfrac23a^{3/2}$; the full tail also carries the $a$-independent
$\beta$-ensemble normalization $c_\beta=\Gamma(\beta/2)/((4\beta)^{\beta/2}2\pi)$ of \eqref{eq:bn}.
There are two distinct asymptotic regimes, and the resurgence question is genuinely different in
each; conflating them is the central pitfall this paper avoids.

\begin{itemize}
\item[\textbf{(L)}] \emph{The level direction.} Fix $\beta$ (the classical point $\beta=2$ is exact,
$\mathcal T_2(a)=1-\exp(-\int_a^\infty(x-a)q^2)$ with $q$ Hastings--McLeod \cite{TW1994}) and let
$a\to\infty$. This expansion is governed by an ODE in the \emph{level} variable $a$ --- Painlev\'e~II
--- and its resurgence is the classical Painlev\'e resurgence: factorially divergent, non-sign-
alternating coefficients, a positive-axis Borel singularity at the instanton action, and a nonzero
Stokes constant $C$ (\S\S\ref{sec:T1proof}--\ref{sec:stokes}).
\item[\textbf{(N)}] \emph{The noise direction.} Fix $a$ and let $\beta\to\infty$. The full $1/\beta$
expansion of the tail --- the dynamical loop factor $F=\sum_n c_n(a)\beta^{-n}$ times the Stirling
expansion of the normalization $c_\beta$ --- solves no ODE in $1/\beta$, so Costin's nonlinear-ODE
theorem does not apply; its Borel-summability is a separate statement, proved by elementary means in
\S\ref{sec:noise}.
\end{itemize}

The bridge between them is the fluctuation operator: the affine Lax map of
Lemma~\ref{lem:lax} shows $\mathcal J$ \emph{is} the PII linearization, which (i) forces the level
expansion (L) to be the Hastings--McLeod trans-series and (ii) supplies the reduced problem whose
Airy structure makes the noise expansion (N) Borel summable.

\section{The Painlev\'e~II backbone (recalled)}\label{sec:backbone}

\begin{proposition}[Weak-noise reduction to PII]\label{prop:pii}
In the large-$\beta$ limit the tail has the large-deviation form $e^{-\beta\Phi(a)}$ with Gaussian
$O(1/\beta)$ fluctuations, obtained by the optimal-fluctuation method on the Riccati-diffusion
representation of $\SAO$; the rate function $\Phi(a)$ solves a Painlev\'e~II equation in the spectral
variable $a$. \tagP\ \emph{(Comtet--Le Doussal--Smith \cite{CLDS}; leading rate and $O(1/\beta)$
cumulants only, with no resurgence analysis.)}
\end{proposition}

\begin{proposition}[Resurgence of the PII backbone]\label{prop:hm}
The Hastings--McLeod solution of PII has a full resurgent trans-series: perturbative and exponential
sectors entwined by resurgent relations, a nonlinear Stokes transition between $s\to\pm\infty$, and,
in the $s\to-\infty$ region, factorially divergent \emph{non-sign-alternating} coefficients forcing
an imaginary trans-series parameter fixed by a median/lateral Borel-resummation reality condition;
integrability is inessential to these properties. \tagP\ \emph{(Cleri--Dunne \cite{CleriDunne}; the
Stokes constant is evaluated in \cite{Dunne2025}, cf.\ \S\ref{sec:stokes}.)}
\end{proposition}

\begin{remark}[Why $q{=}1$ is provable and the cusp $q{=}2$ is not]\label{rem:crutch}
The $q{=}1$ tail reduces to an integrable ODE backbone (PII) already known resurgent
(Prop.~\ref{prop:hm}); the cusp law $\mathcal W_\beta$ has no Painlev\'e reduction (isomonodromy
fixed point), so it is established only numerically. This paper is the theorem-grade sibling; the
``integrability inessential'' clause is the bridge to the cusp case.
\end{remark}

\section{Main results}\label{sec:results}

\begin{theorem}[T1: the level expansion is the Hastings--McLeod PII trans-series]\label{thm:T1}
Under the affine map of Lemma~\ref{lem:lax} the tail's fluctuation operator is the exact Painlev\'e~II
linearization, and the instanton is the $\alpha=\tfrac12$ Airy special solution of PII. Consequently
the level expansion \textbf{(L)} of \S\ref{sec:two} is the Hastings--McLeod trans-series: an ODE in
$a$ of rank-one/irregular type at $a=\infty$, so by Costin's theorem \cite{Costin} it is Borel
summable (median/lateral), with Borel singularities at the instanton action and the
non-sign-alternating reality structure of Prop.~\ref{prop:hm}; its median sum is real and, at
$\beta=2$, equals $\mathcal T_2(a)$. \tagPr\ \emph{(modulo the cited PII/Costin machinery; \S\ref{sec:T1proof}.)}
\end{theorem}

\begin{theorem}[T2: Stokes data --- location, shift, and the exact Stokes constant]\label{thm:T2}
\emph{(a) Location-fixing.} The Borel-singularity locations of \textbf{(L)} are noise-independent
(equal to the deterministic instanton action), the noise average shifting only the connection/Stokes
data --- the connection root $\Lambda_0$ (whose shift is (b)) and the Stokes constant (c).
\emph{(b) Connection-root shift.} $\SAO$ is self-adjoint, so its spectrum is real for every
realization; $\E[\Lambda_0]=\Lambda_0^{\det}+\beta^{-1}\Omega_1+O(\beta^{-2})$ is real with
$\Omega_1=-1.12\pm0.01$ a finite renormalized Airy-Green's chaos integral.
\emph{(c) The Stokes constant.} The $s\to-\infty$ Hastings--McLeod coefficients are
non-sign-alternating; the Borel singularity sits on the positive real axis at action
$A=\tfrac{2\sqrt2}{3}$, index $\gamma=-\tfrac12$, and the Stokes constant is
\[
C=\sqrt{\tfrac{2}{3\pi^3}}=\tfrac1\pi\sqrt{\tfrac{2}{3\pi}}=0.146632271193848478\ldots,
\]
computed here to $28$ digits from an exact rational recursion and matching the published value.
\emph{(d) The $a$-dependence is algebraic.} $C$ factorizes as $\tfrac1\pi\cdot\sqrt{2/(3\pi)}$, where
$\tfrac1\pi$ is the transverse Kramers determinant and $\sqrt{2/(3\pi)}=\lVert\psi_0\rVert/\sqrt{2\pi}$
is the Gelfand--Yaglom zero-mode Jacobian of the tanh-kink escape instanton; the whole $a$-dependence
of the tail-specific Stokes constant is the elementary scaling factor of this reduction, carrying no
new transcendental. \tagPr\ \emph{(\S\S\ref{sec:T2proof}, \ref{sec:stokes}.)}
\end{theorem}

\begin{theorem}[T3: the noise expansion is Borel summable]\label{thm:T3}
The noise expansion \textbf{(N)} of the tail is Borel summable along the noise ray in the median
sense --- a two-sector resurgent trans-series. Its \emph{dynamical escape} sector (the loop factor
$F=\sum_n c_n(a)\beta^{-n}$) has a positive-real-axis Borel singularity, its median ambiguity being
the two-instanton term $e^{-2\beta\Phi(a)}$, and is median summable; its \emph{Coulomb-gas} sector (the $\Gamma(\beta/2)$ normalization
in \eqref{eq:bn}) has alternating-factorial coefficients and imaginary-axis Borel singularities, and
is ordinary Borel summable. The dynamical sector is proved elementarily via an Airy-type integral
representation of the reduced escape rate (no nonlinear-ODE theorem required).
\tagPr\ \emph{(scaling regime; \S\ref{sec:noise}.)}
\end{theorem}

\begin{conjecture}[Far frontier]\label{conj:frontier}
A direct exact-WKB / Voros-symbol resurgence for the white-noise operator $\SAO$ itself, of which
Theorems~\ref{thm:T1}--\ref{thm:T3} are the PII-mediated shadow; no exact-WKB framework for a
random-potential Schr\"odinger operator currently exists. \tagX
\end{conjecture}

\section{Proof of Theorem~\ref{thm:T1}: the Lax map and the level resurgence}\label{sec:T1proof}

\subsection{Exact reduction.}
With the Riccati variable $p=\psi'/\psi$ for $-\psi''+(x+2\eps b')\psi=\Lambda\psi$ at level
$\Lambda=-a$, one has the It\^o SDE $\dd p=((x+a)-p^2)\dd x+2\eps\dd B$, $p(0^+)=+\infty$, and by
Sturm oscillation $\mathcal T_\beta(a)=\Prob(\text{explosion to }-\infty)$. The no-explosion
probability $u(p,x)$ solves the \emph{exact} backward-Kolmogorov equation
\begin{equation}\label{eq:bk}
\partial_x u+((x+a)-p^2)\partial_p u+2\eps^2\partial_p^2 u=0.
\end{equation}
\tagPr

\subsection{Weak-noise WKB and the instanton.}
Insert $u=e^{-W/\eps^2}(A_0+\eps^2 A_1+\cdots)$. The eikonal order gives the Hamilton--Jacobi
equation with Freidlin--Wentzell Hamiltonian $H=\theta b+2\theta^2$, $b=(x+a)-p^2$; the
zero-energy fluctuation branch $\theta=-b/2$ yields the instanton
\begin{equation}\label{eq:inst}
\dot\pinst=\pinst^2-(x+a),\qquad \pinst=-\varphi'/\varphi,\quad \varphi''=(x+a)\varphi\ (\text{Airy}),
\end{equation}
with action $\Phi(a)=\tfrac12\int b^2\dd x=\tfrac23 a^{3/2}$ at leading order, matching
Prop.~\ref{prop:anchor}. Collecting the orders $\eps^{2n}$ along the instanton characteristics gives
the linear transport hierarchy $\mathcal L_0 A_n=\kappa A_n-2\partial_p^2 A_{n-1}$ that generates the
$c_n(a)=A_n(+\infty,0)$; each $c_n$ is a finite iterated Airy integral (after removal of the
translation zero mode \cite{SGG}, which leaves the finite $a^{-3\beta/4}$ prefactor). \tagPr

\subsection{The fluctuation operator.}
The Hessian of the FW action $I=\tfrac18\int(\dot p-b)^2$ about \eqref{eq:inst} is
$\delta^2 I=\tfrac14\int\eta\,\mathcal J\,\eta$ with
\begin{equation}\label{eq:jacobi}
\mathcal J=-\partial_x^2+V_{\mathcal J}(x),\qquad V_{\mathcal J}(x)=6\pinst^2-2(x+a),
\end{equation}
by integration by parts using $\dot\pinst=\pinst^2-(x+a)$. The coefficient $6$ on $\pinst^2$ is the
Painlev\'e~II fingerprint. \tagPr

\subsection{The Lax-pair map (the crux).}
\begin{lemma}[Instanton $=$ Airy solution of PII${}_{\alpha=1/2}$; $\mathcal J=$ its linearization]
\label{lem:lax}
Set $\mu:=-2^{1/3}$ ($\mu^3=-2$), $s:=\mu(x+a)$, $q(s):=\mu^{-1}\pinst$. Then $q$ solves the Riccati
equation $q'=q^2+s/2$, hence Painlev\'e~II with parameter $\alpha=\tfrac12$, $q''=2q^3+sq+\tfrac12$,
and under the same change of variables
\[
\mathcal J=-\partial_x^2+6\pinst^2-2(x+a)=\mu^2\bigl(-\partial_s^2+6q^2+s\bigr)=\mu^2\,\mathcal L_{\mathrm{PII}},
\]
the exact linearization of PII about $q$. \tagPr\ \emph{(algebra verified symbolically and numerically
to machine precision.)}
\end{lemma}
\begin{proof}
$q_s=\mu^{-2}\pinst'=\mu^{-2}(\pinst^2-z)$ with $z=x+a$; since $q^2=\mu^{-2}\pinst^2$ and
$\mu^{-2}z=\mu^{-3}s=-\tfrac12 s$, $q_s=q^2+\tfrac12 s$, and differentiating gives PII with
$\alpha=\tfrac12$. For $\mathcal J$: $-\partial_x^2=-\mu^2\partial_s^2$ and
$6\pinst^2-2z=\mu^2(6q^2)-2\mu^{-1}s=\mu^2(6q^2+s)$ (using $-2\mu^{-1}=\mu^2$). \qed
\end{proof}
Equation $q'=q^2+s/2$ is the defining Riccati of the classical \emph{Airy special solutions} of PII
($q=-w'/w$, $w''=-\tfrac{s}{2}w$), which exist at $\alpha\in\tfrac12+\Z$ \cite{FIKN,Clarkson}; the
instanton \eqref{eq:inst} is exactly this object. The Flaschka--Newell linear problem \cite{FN1980}
restricted to the Airy locus, with its Jimbo--Miwa--Ueno $\tau$-function \cite{JMU1981}, supplies the
isomonodromic frame; at $\beta=2$ the tail is exactly the Hastings--McLeod expression
$\mathcal T_2(a)=1-\exp(-\int_a^\infty(x-a)q^2\dd x)$ \cite{TW1994}, fixing the normalization.

\subsection{Closing the level direction.}
By Lemma~\ref{lem:lax} the transport hierarchy is the linearized-PII recursion, so the level
expansion \textbf{(L)} is the Hastings--McLeod trans-series. Painlev\'e~II $q''=2q^3+sq+\alpha$ has a
rank-one irregular singularity at $s=\infty$ for \emph{every} $\alpha$, so Costin's theorem
\cite{Costin} gives Borel summability in the generalized median-path sense with a one-to-one,
Stokes-jumping trans-series$\leftrightarrow$solution correspondence; the Stokes constant is
non-vanishing \cite{CCK}. This proves Theorem~\ref{thm:T1}. The prefactor $a^{-3\beta/4}$ is
\emph{eikonal} --- its $\log a$ coefficient in $-\log\mathcal T_\beta$ is $\tfrac34\beta$, part of the
leading $\beta$-exponential --- so the PII parameter is $\beta$-independent and does not enter the summability class.
\tagPr\,/\,\tagP

\begin{remark}[The $\alpha=\tfrac12$ vs $\alpha=0$ family --- a sharpening]\label{rem:family}
The leading tail \emph{saddle} is the $\alpha=\tfrac12$ Airy solution (Lemma~\ref{lem:lax}); the full
$\beta=2$ \emph{distribution} is the $\alpha=0$ Hastings--McLeod transcendent --- the ordinary
saddle-vs-distribution distinction, not a $\beta$-dependence. Since Costin's class contains PII for
all $\alpha$, Theorem~\ref{thm:T1} holds regardless; pinning the correspondence sharpens \emph{which}
PII, not \emph{whether}. \tagT
\end{remark}

\section{The Stokes constant and its algebraic $a$-dependence}\label{sec:stokes}

This section proves Theorem~\ref{thm:T2}(c,d).

\subsection{Exact value of $C$.}\label{ssec:C}
Substituting $q(s)=\sqrt{-s/2}\sum_{k\ge0}a_k(-s)^{-3k}$ into $q''=2q^3+sq$ (the $s\to-\infty$
Hastings--McLeod branch) and matching powers gives the closed-form rational recursion
\begin{equation}\label{eq:rec}
a_m=\tfrac12\Big[a_{m-1}\,p_{m-1}(p_{m-1}-1)-\!\!\!\sum_{\substack{i+j+l=m\\ 0\le i,j,l\le m-1}}\!\!\! a_i a_j a_l\Big],
\quad p_{m-1}=\tfrac12-3(m-1),\quad a_0=1,
\end{equation}
so $a_1=-\tfrac18$, $a_2=-\tfrac{73}{128}$, $a_3=-\tfrac{10657}{1024},\dots$ (the known
Hastings--McLeod values). All $a_{k\ge1}<0$ (non-sign-alternating), so the Borel singularity is on
the positive real axis and the physical asymptotics is the median of the two lateral Borel sums. The
growth is $a_k\sim -C\,\Gamma(2k+\gamma)(9/8)^k$; carrying \eqref{eq:rec} to $120$ orders in exact
arithmetic and Richardson-extrapolating gives
\[
\gamma=-\tfrac12\ \text{(to $25$ digits)},\qquad
C=0.146632271193848478\ldots\ \text{(stable to $28$ digits)},
\]
whence the action $A=\tfrac{2\sqrt2}{3}$ ($9/8=1/A^2$). The value is identified in closed form,
\begin{equation}\label{eq:C}
\boxed{\,C=\sqrt{\tfrac{2}{3\pi^3}}=\tfrac1\pi\sqrt{\tfrac{2}{3\pi}}\,,}
\end{equation}
agreeing with the computed number to $28$ digits and with the published Hastings--McLeod Stokes
constant of \cite{Dunne2025} (eq.~(2.16), which writes it in exactly the factored form
$\tfrac1\pi\sqrt{2/(3\pi)}$; cf.\ \cite{KapaevHM,BothnerRHP}). Our coefficients reproduce not only the
leading constant but the first subleading term of \cite{Dunne2025}~eq.~(2.16): forming
$R_n=a_n/[-C\,\Gamma(2n-\tfrac12)(9/8)^n]$, the quantity $(R_n-1)(2n-\tfrac32)$ Richardson-extrapolates
to $-\tfrac{17}{72}$ to $16$ digits. \emph{Equation \eqref{eq:C} refutes the natural guess
$C\stackrel?=1/(4\sqrt\pi)=0.14105$, which is $3.8\%$ too small.} \tagPr

Two independent checks confirm \eqref{eq:C}. \emph{(i) Borel-space.} With $c_k=a_k/(2k)!$ the
transform $\sum_k a_k\xi^{2k}/(2k)!$ has a bounded $(A^2-\xi^2)^{1/2}$ branch at $\xi=A$; the branch
exponent extracts to $-\tfrac12$ and the amplitude to $C$. \emph{(ii) ODE.} Direct high-precision
integration of the true Hastings--McLeod transcendent shows its optimally truncated series remainder
decaying at exactly the rate $e^{-A(-s)^{3/2}}$, the action of \eqref{eq:C}. \tagPr

\subsection{The $a$-dependence is a zero-mode Jacobian.}\label{ssec:adep}
Under $p=\sqrt a\,P$ the backward-Kolmogorov PDE \eqref{eq:bk} becomes (dividing through by $a$)
\begin{equation}\label{eq:scaled}
a^{-3/2}\partial_X u+(X-P^2)\partial_P u+2\eps^2 a^{-3/2}\partial_P^2 u=0,\qquad X:=(x+a)/a,
\end{equation}
verified symbolically: $a$ enters only through the single combination $g:=\beta a^{3/2}$ (the
effective inverse-noise $\eps_{\mathrm{eff}}^2=1/g$). The dominant boundary escape thus reduces to an
$a$-independent cubic-barrier problem, with escape rate
\begin{equation}\label{eq:R}
R(g)=\tfrac1\pi\,e^{-\frac23 g}\,\tilde Q(2/g)^{-2},\qquad
\tilde Q(D)=\sum_{j\ge0}\tilde q_j D^j,\quad \tilde q_j=\frac{(6j)!}{576^{\,j}\,(3j)!\,(2j)!},
\end{equation}
(the two Laplace peaks of the mean-first-passage integral decouple to all orders). The reduced
fluctuation series $\tilde Q(2/g)^{-2}=\sum_m d_m g^{-m}$ has, certified to $25$ digits,
$d_m\sim-\tfrac1\pi\Gamma(m)(3/2)^m$: index $\gamma'=0$ and \emph{reduced Stokes constant}
$S_0=1/\pi$. Comparing to \S\ref{ssec:C}, the frozen datum $(S_0,\gamma')=(1/\pi,0)$ is promoted to
$(C,\gamma)=(C,-\tfrac12)$ by a \emph{single joint shift}: index $-\tfrac12$ and constant factor
$C/S_0=\pi C=\sqrt{2/(3\pi)}$ (agreement $1.4\times10^{-9}$ from the matched-order ratio). This factor
is a Gelfand--Yaglom zero-mode Jacobian.

\begin{proposition}[The $x$-extension factor is $\lVert\psi_0\rVert/\sqrt{2\pi}$]\label{prop:gy}
The reduced escape instanton is the tanh kink $P_\star(\tau)=-\tanh\tau$ (the uphill solution
$\dot P=V'(P)=P^2-1$ of the cubic barrier $V=\tfrac13P^3-P$, well $P{=}{+}1$ to barrier $P{=}{-}1$).
Its translation zero mode $\psi_0=\dot P_\star=-\sech^2\tau$ has
\begin{equation}\label{eq:virial}
\lVert\psi_0\rVert^2=\int_\R\sech^4\tau\,\dd\tau=\tfrac43=\Delta V=2A_g,\qquad A_g:=\tfrac23,
\end{equation}
by the virial identity $\int\dot P^2\dd\tau=\int V'\dd P=\Delta V$. Hence the zero-mode
(collective-coordinate) Jacobian is $\sqrt{\lVert\psi_0\rVert^2/2\pi}=\sqrt{2/(3\pi)}=\sqrt{A_g/\pi}$,
and $C=S_0\cdot\sqrt{\lVert\psi_0\rVert^2/2\pi}=\tfrac1\pi\sqrt{2/(3\pi)}=\sqrt{2/(3\pi^3)}$ --- an
exact algebraic identity, its inputs $S_0=1/\pi$ and $\lVert\psi_0\rVert^2=4/3$ each certified above.
The accompanying half-power $\sqrt g$ is realized as the $-\tfrac12$ index shift (one continuous zero
mode). \tagPr
\end{proposition}

Thus the $a$-dependence of the tail-specific Stokes constant is purely algebraic (the elementary
scaling of \eqref{eq:scaled}), the transcendental content being the $a$-independent $C$; and $C$
itself decomposes into a transverse Kramers determinant $1/\pi$ and the $x$-translation zero-mode
Jacobian $\sqrt{2/(3\pi)}$ --- the very mode whose removal produces the eikonal prefactor
$a^{-3\beta/4}$ (\S\ref{sec:T1proof}). This proves Theorem~\ref{thm:T2}(c,d).

\section{Proof of Theorem~\ref{thm:T3}: Borel-summability along the noise ray}\label{sec:noise}

Unlike the level direction, the noise loop expansion solves no ODE in $1/\beta$; we prove its
Borel-summability directly, and elementarily, from the reduced problem of \S\ref{ssec:adep}.

\subsection{The dynamical escape sector.}
By \eqref{eq:R}, $F(g)=\tilde Q(2/g)^{-2}=\sum_m d_m g^{-m}=\pi e^{\frac23 g}R(g)$ is the reduced
fluctuation series; at fixed $a$ the noise series is $d_m a^{-3m/2}\beta^{-m}$ (so
$c_n(a)=d_n a^{-3n/2}$ and $F(a;\beta)=F(g)$ with $g=\beta a^{3/2}$ in the scaling regime),
non-sign-alternating (all $d_m<0$), so its Borel singularity is on the positive real (noise) ray.

\begin{lemma}[Airy integral representation]\label{lem:airy}
$\tilde Q(D)=Q(D)/\sqrt\pi$ with the cubic-phase integral
$Q(D)=\int_\Gamma e^{-t^2+(\sqrt D/3)t^3}\dd t$, $\Gamma$ the steepest-descent contour through the
saddle $t=0$; the affine map $t=D^{-1/6}s+D^{-1/2}$ sends
$Q(D)=e^{-2/(3D)}D^{-1/6}\int e^{s^3/3-\zeta s}\dd s$, $\zeta=D^{-2/3}$ --- a genuine Airy integral. Its
asymptotic expansion as $D\to0^+$ is $\sqrt\pi\,\tilde Q(D)$.
\end{lemma}
\emph{Numerical certification.} $Q(D)/\sqrt\pi$ against the optimally truncated $\tilde Q(D)$ series:
relative difference $2.2\times10^{-6}$ at $D=0.12$, improving to $2.1\times10^{-11}$ at $D=0.06$ (the
signature of a correct asymptotic-to-integral match). \tagPr

\begin{proposition}[Borel-transform structure]\label{prop:borel}
The Borel transform $\mathcal B(\xi)=\sum_j \tilde q_j\,\xi^j/j!$ is analytic in $\C\setminus[4/3,\infty)$
with a single logarithmic branch point at $\xi=4/3=\Delta V$ (the cubic-barrier action; it generates
the exponential $e^{-(4/3)/D}=e^{-\beta\Phi}$) and subexponential growth off the cut.
\end{proposition}
\emph{Certification.} (a) Pringsheim: all Borel coefficients $\tilde q_j/j!>0$, so the dominant
singularity is on $\R_+$. (b) Domb--Sykes: $(\tilde q_j/j!)/(\tilde q_{j-1}/(j-1)!)\to 3/4$, radius
$4/3$. (c) Amplitude: $j(4/3)^j(\tilde q_j/j!)\to 1/(2\pi)$, i.e.\ $\mathcal B\sim-(2\pi)^{-1}\log(1-3\xi/4)$
(logarithmic branch). (d) Borel--Pad\'e: the genuine poles all lie at $\mathrm{Re}\,\xi\ge1.344\approx4/3$
(Froissart doublets excluded via numerator-zero pairing), tracing the single cut. \tagPr

\begin{proof}[Proof of the dynamical sector of Theorem~\ref{thm:T3}]
$\tilde q_j\sim(2\pi)^{-1}\Gamma(j)(3/4)^j$ (Stirling), so $\tilde Q$ is Gevrey-1 and, by
Lemma~\ref{lem:airy}, is the asymptotic series of the convergent Airy integral $Q(D)/\sqrt\pi$. By
Prop.~\ref{prop:borel} the only Borel singularity lies on the summation ray $\R_+$; the two lateral
Laplace transforms $\mathcal L_{\pm\theta}\mathcal B$ ($\theta\to0^+$) therefore exist (analytic +
subexponential off the cut) and are complex conjugates (real $\tilde q_j$), so their \emph{median}
$\tfrac12(\mathcal L_{+\theta}+\mathcal L_{-\theta})$ is real, angle-independent, and equals
$Q(D)/\sqrt\pi$ (Nevanlinna's theorem for the Gevrey-1 remainder bounds \cite{Sokal}, together with
\'Ecalle/Costin median summation \cite{Costin} for the on-axis singularity). The lateral ambiguity
$i\,e^{-(4/3)/D}$ is the one-instanton Stokes term. Finally, median-Borel-summable series form a ring
closed under inversion at units; since $\tilde Q(0)=1$, the reciprocal square $F=\tilde Q^{-2}$ is
median Borel summable, with median sum $\pi e^{\frac23 g}R(g)$. A direct median Borel--Laplace of the
$d_m$ reconstructs the exact escape rate ($F$ at $g=9$) to $1.2\times10^{-3}$. \qed
\end{proof}

\subsection{The Coulomb-gas sector.}
The $\Gamma(\beta/2)$ normalization in \eqref{eq:bn} contributes, through
$\log\Gamma(\beta/2)$, the Stirling series $\sum_k \tfrac{B_{2k}}{2k(2k-1)}(2/\beta)^{2k-1}$: its
coefficient of $\beta^{-(2k-1)}$ is $B_{2k}\,2^{2k-1}/(2k(2k-1))\sim(-1)^{k+1}(2k-2)!/\pi^{2k}$,
\emph{alternating} factorial (verified: the ratio $\to 1/\pi^2$ to $6$ digits), with Borel transform
$\sim(1/\pi)\arctan(\xi/\pi)$ and branch points at $\xi=\pm i\pi$ on the imaginary axis. This is the
classical Borel-summable (ordinary, no median) series --- the $\Gamma$-function / Coulomb-gas
instantons $e^{-2\pi i n\beta/2}$. \tagPr\,/\,\tagP

\begin{remark}[Elementary, and honestly scoped]\label{rem:elem}
The level direction (L) genuinely needs Costin's nonlinear PII theorem; the noise ray (N) does not,
because the reduced problem is the \emph{linear} Airy structure of Lemma~\ref{lem:airy}, whence
Watson/Nevanlinna \cite{Sokal} plus median summation suffice. This proves the dynamical sector in the
scaling regime $g=\beta a^{3/2}\to\infty$ (where the frozen escape is exact) and the Coulomb-gas
sector exactly; the uniform statement at fixed finite $a$ adds the $O(a^{-3/2})$ $x$-extension
corrections of \S\ref{ssec:adep} and would combine this proof with the T1 PII map. \tagT
\end{remark}

\section{Location and shift (Theorem~\ref{thm:T2}(a,b))}\label{sec:T2proof}

\subsection{Location-fixing (a).}
The Borel-singularity location is the WKB action of the noise-dressed problem,
$A[b']=\int\sqrt{(x+a)+2\eps b'}\dd x$. Its first-order shift $A_1=\int b'/\sqrt{x+a}$ has
$\E[A_1]=0$ (the instanton is a saddle), fixing the location at $O(\eps)$; at $O(\eps^2)$,
$\E[A_2]=-\tfrac12\delta(0)\int(x+a)^{-3/2}$ is a coincident-point self-contraction, hence
$\delta(0)$-divergent --- a formal manipulation whose rigorous content is that the finite noise
information passes to the Stokes/connection data. The
periods are Voros integrals over the noise-independent Airy spectral curve $y^2=x+a$; the noise,
subleading at the rank-$3/2$ irregular singularity, preserves the formal-monodromy exponents to all
orders (isomonodromy rigidity), moving only the Stokes matrices. \tagPr\ \emph{(both orders verified
numerically; rough-potential rigor of the rigidity is a remaining technical point.)}

\subsection{The connection-root shift $\Omega_1$ (b).}
$\SAO$ is self-adjoint, so its spectrum is real for \emph{every} realization (verified
$\max|\mathrm{Im}\,\Lambda_n|=0$). Second-order perturbation theory with $W=2\eps b'$ gives
\begin{equation}\label{eq:om1}
\E[\Lambda_0]=\Lambda_0^{\det}+\eps^2\Omega_1+O(\eps^4),\qquad
\Omega_1=4\!\sum_{n\ge1}\frac{\langle\psi_0^2|\psi_n^2\rangle}{\Lambda_0-\Lambda_n}
=-4\!\int_0^\infty\!\psi_0^2\,G_{\mathrm{red}}(x,x)\dd x,
\end{equation}
finite because the Airy Green's function spreads the $\delta(0)$ contraction (summand $\sim n^{-4/3}$).
Two independent evaluations agree: $-1.121$ (sum-over-states) and $-1.128\pm0.009$
(direct-diagonalization Monte Carlo with a Malliavin/Stein variance-reduced estimator \cite{Viens},
$\eps$-independent), so $\Omega_1=-1.12\pm0.01$ and
$\E[\Lambda_0]=2.3381-1.12/\beta+O(\beta^{-2})$. Being real (self-adjointness) is the structural
contrast with the cusp, whose unbounded-below Weber operator has a complex resonance
$\lambda_0=0.8896(1-i)$ and complex $\Omega_2=-0.451+0.352i$. \tagPr

\section{Honest ledger and positioning}\label{sec:ledger}

\emph{Established here (modulo cited machinery).} \textbf{T1} --- the level expansion is the
Hastings--McLeod PII trans-series --- via the exact reduction \eqref{eq:bk}, the instanton
\eqref{eq:inst}, the Jacobi operator \eqref{eq:jacobi}, and the Lax map (Lemma~\ref{lem:lax}),
whereupon Costin \cite{Costin} closes summability. \textbf{T2} --- location-fixing (a), the real shift
$\Omega_1=-1.12$ (b, two evaluations), the exact Stokes constant $C=\sqrt{2/(3\pi^3)}$ (c, $28$
digits, literature-matched), and its algebraic $a$-dependence as the
tanh-kink zero-mode Jacobian (d, Prop.~\ref{prop:gy}, an exact identity). \textbf{T3} --- the noise loop
expansion is median Borel summable, a two-sector trans-series proved elementarily via the Airy
integral (Lemma~\ref{lem:airy}, Prop.~\ref{prop:borel}). \emph{Cited (turnkey):}
\cite{RRV,DV,BorotNadal,CLDS,CleriDunne,Costin,CCK,SGG,TW1994,FN1980,JMU1981,FIKN,Clarkson,Dunne2025,KapaevHM,BothnerRHP,Sokal,Viens}.
\emph{Refinements:} the $\alpha=\tfrac12/\alpha=0$ family assignment (Rem.~\ref{rem:family}); uniform
fixed-$a$ noise summability and rough-potential rigidity rigor (Rem.~\ref{rem:elem}). \emph{Open
frontier:} Conjecture~\ref{conj:frontier}.

\paragraph{Positioning.} This is, to our knowledge, the first Borel-summable resurgent trans-series
for a random-operator edge law, and it is resurgent in two directions at once: the Painlev\'e~II
level resurgence (Stokes constant $C=\sqrt{2/(3\pi^3)}$) and the weak-noise loop resurgence (two
sectors, dynamical escape and Coulomb-gas). It is the theorem-grade sibling of the cusp companion,
where the absence of a Painlev\'e reduction forces a numerical partial resurgent representation; here
the Airy/PII backbone makes the analogous statements theorems, and a clean structural dichotomy
emerges --- the cusp's complexity is a \emph{resonance} (non-self-adjoint, unbounded-below)
phenomenon, whereas $\TW$'s is a \emph{trans-series} (self-adjoint, median-resummation) one. The one
genuinely open piece common to both is Conjecture~\ref{conj:frontier}: a direct exact-WKB of the
random operator.

\begin{thebibliography}{99}

\bibitem{RRV} J.~Ram\'irez, B.~Rider, B.~Vir\'ag, \emph{Beta ensembles, stochastic Airy spectrum,
and a diffusion}, J. Amer. Math. Soc. \textbf{24} (2011) 919--944; arXiv:math/0607331.

\bibitem{DV} L.~Dumaz, B.~Vir\'ag, \emph{The right tail exponent of the Tracy--Widom-$\beta$
distribution}, Ann. Inst. H. Poincar\'e Probab. Statist. \textbf{49} (2013) 915--933;
arXiv:1102.4818.

\bibitem{BorotNadal} G.~Borot, C.~Nadal, \emph{Right tail asymptotic expansion of Tracy--Widom beta
laws}, Random Matrices Theory Appl. \textbf{1} (2012) 1250006; arXiv:1111.2761.

\bibitem{CLDS} A.~Comtet, P.~Le Doussal, N.~R.~Smith, \emph{The Tracy--Widom distribution at large
Dyson index}, J. Stat. Phys. \textbf{193} (2026) Art.~46; arXiv:2510.14433.

\bibitem{CleriDunne} N.~J.~Cleri, G.~V.~Dunne, \emph{Resurgent trans-series for generalized
Hastings--McLeod solutions}, J. Phys. A: Math. Theor. \textbf{53} (2020) 305201; arXiv:2002.06270.

\bibitem{Costin} O.~Costin, \emph{On Borel summation and Stokes phenomena of nonlinear
differential systems}, Duke Math. J. \textbf{93} (1998) 289--344; arXiv:math/0608408.

\bibitem{CCK} O.~Costin, R.~D.~Costin, M.~Kohut, \emph{Rigorous bounds of Stokes constants for some
nonlinear ODEs at rank-one irregular singularities}, Proc. R. Soc. A \textbf{460} (2004)
3631--3641; arXiv:math/0608316.

\bibitem{SGG} T.~Schorlepp, T.~Grafke, R.~Grauer, \emph{Symmetries and zero modes in sample path
large deviations}, J. Stat. Phys. \textbf{190} (2023) 50; arXiv:2208.08413.

\bibitem{Sokal} A.~D.~Sokal, \emph{An improvement of Watson's theorem on Borel summability},
J. Math. Phys. \textbf{21} (1980) 261--263.

\bibitem{Viens} F.~G.~Viens, \emph{Stein's lemma, Malliavin calculus, and tail bounds, with
application to polymer fluctuation exponent}, Stochastic Process. Appl. \textbf{119} (2009)
3671--3698; arXiv:0901.0383.

\bibitem{TW1994} C.~A.~Tracy, H.~Widom, \emph{Level-spacing distributions and the Airy kernel},
Comm. Math. Phys. \textbf{159} (1994) 151--174.

\bibitem{FN1980} H.~Flaschka, A.~C.~Newell, \emph{Monodromy- and spectrum-preserving deformations.
I}, Comm. Math. Phys. \textbf{76} (1980) 65--116.

\bibitem{JMU1981} M.~Jimbo, T.~Miwa, K.~Ueno, \emph{Monodromy preserving deformation of linear
ordinary differential equations with rational coefficients. I. General theory and $\tau$-function},
Physica D \textbf{2} (1981) 306--352.

\bibitem{FIKN} A.~S.~Fokas, A.~R.~Its, A.~A.~Kapaev, V.~Yu.~Novokshenov, \emph{Painlev\'e
Transcendents: The Riemann--Hilbert Approach}, Math. Surveys Monogr. \textbf{128}, AMS, 2006.

\bibitem{Clarkson} P.~A.~Clarkson, \emph{Painlev\'e equations --- nonlinear special functions},
J. Comput. Appl. Math. \textbf{153} (2003) 127--140.

\bibitem{Dunne2025} G.~V.~Dunne, \emph{Introductory lectures on resurgence}, arXiv:2511.15528
(2025). [Hastings--McLeod Stokes constant $\tfrac1\pi\sqrt{2/(3\pi)}$, eq.~(2.16).]

\bibitem{KapaevHM} A.~A.~Kapaev, \emph{Quasi-linear Stokes phenomenon for the Hastings--McLeod
solution of the second Painlev\'e equation}, arXiv:nlin/0411009 (2004).

\bibitem{BothnerRHP} T.~Bothner, \emph{On the origins of Riemann--Hilbert problems in mathematics},
Nonlinearity \textbf{34} (2021) R1--R73; arXiv:2003.14374.

\end{thebibliography}

\end{document}
